% !TEX program = xelatex
\documentclass{英语试卷}
\begin{document}
\examtitle{炎德·英才大联考雅礼中学 2026 届高三月考试卷(一)}{英语}
\par 炎德·英才大联考雅礼中学 2026 届高三月考试卷(一)\par
\par 高三英语 试题\par
\begin{notice}
\par 1. 答卷前，考生务必将自己的姓名、准考证号填写在答题卡上。\par
\par 2. 回答选择题时，选出每小题答案后，用铅笔把答题卡上对应题目的答案标号涂黑。如需改动，用橡皮擦干净后，再选涂其他答案标号。回答非选择题时，将答案写在答题卡上，写在本试卷上无效，\par
\par 3. 考试结束后，将本试卷和答题卡一并交回。\par
\end{notice}
\parttitle{第一部分 听力（共两节，满分 30 分）}
\par 做题时，先将答案标在试卷上，录音内容结束后，你将有两分钟的时间将试卷上的答案转涂到答题卡上。\par
\sectitle{第一节（共5小题：每小题1.5分，满分7.5分）}
\par 听下面 5 段对话。每段对话后有一个小题，从题中所给的 A、B、C 三个选项中选出最佳选项。听完每段对话后，你都有 10 秒钟的时间来回答有关小题和阅读下一小题。每段对话仅读一遍。\par
\question{1. What kind of movies does the woman like best?}
\begin{choices}(3)
  \item Action movies.
  \item Musicals.
  \item Science fiction
\end{choices}
\question{2. What was the weather like during the vacation?}
\begin{choices}(3)
  \item Rainy.
  \item Cool.
  \item Too hot.
\end{choices}
\question{3. What time is it probably now?}
\begin{choices}(3)
  \item About 1: 00 p.m.
  \item About 2: 00 p.m.
  \item About 3: 00 p.m.
\end{choices}
\question{4. What event are the speakers talking about?}
\begin{choices}(3)
  \item A flower show.
  \item An art show.
  \item A concert.
\end{choices}
\question{5. What attracts the woman most about the phone?}
\begin{choices}(3)
  \item The camera.
  \item The touchscreen.
  \item The GPS system.
\end{choices}
\sectitle{第二节（共15小题：每小题1.5分，满分22.5分）}
\par 听下面 5 段对话或独白。每段对话或独白后有几个小题，从题中所给的 A、B、C 三个选项中选出最佳选项。听每段对话或独白前，你将有时间阅读各个小题，每小题 5 秒钟；听完后，各小题将给出 5 秒钟的作答时间。每段对话或独白读两遍。 听第 6 段材料，回答第 6、7 题。\par
\question{6. Why is Harry unwilling to join the woman?}
\begin{choices}(3)
  \item He wants to watch TV.
  \item He has a pain in his knee.
  \item He wishes to play with Helen.
\end{choices}
\question{7. What will the woman probably do next?}
\begin{choices}(3)
  \item Stay at home.
  \item Do some exercise.
  \item Take Harry to hospital.
\end{choices}
\par 听第 7 段材料，回答第 8 至 10 题。\par
\question{8. Why did the man lose control of the car?}
\begin{choices}(3)
  \item Because he had a bet with others.
  \item Because he probably hit an icy patch.
  \item Because he wanted to drive around the corner.
\end{choices}
\question{9. How was the driver of the other car?}
\begin{choices}(3)
  \item She was nothing serious.
  \item She was frightened at the news.
  \item She was carried away in an ambulance.
\end{choices}
\question{10. What did the man probably want to do at last?}
\begin{choices}(3)
  \item Call the police.
  \item Go to hospital.
  \item Call his family.
\end{choices}
\par 听第 8 段材料，回答第 11 至 13 题，\par
\question{11. Where does the conversation probably take place?}
\begin{choices}(3)
  \item In an office.
  \item At home.
  \item At a restaurant.
\end{choices}
\question{12. What will the speakers do tomorrow evening?}
\begin{choices}(3)
  \item Go to a concert.
  \item Visit a friend.
  \item Work extra hours.
\end{choices}
\question{13. Who is Alice going to call?}
\begin{choices}(3)
  \item Mike.
  \item Joan.
  \item Catherine.
\end{choices}
\par 听第 9 段材料，回答第 14 至 17 题。\par
\question{14. What does the competition try to find this year?}
\begin{choices}(3)
  \item The best dancer.
  \item The best actor.
  \item The best singer.
\end{choices}
\question{15. How should people enter the competition?}
\begin{choices}(3)
  \item By fax.
  \item By phone.
  \item Through the Internet.
\end{choices}
\question{16. What time does the sports hall close this week?}
\begin{choices}(3)
  \item At 10:30 p.m.
  \item At 9:00 p.m.
  \item At 8:00 p.m.
\end{choices}
\question{17. What does the woman say about the new swimming pool?}
\begin{choices}(3)
  \item It will be free for the first week.
  \item It needs to be booked in advance.
  \item It is much bigger than the old one.
\end{choices}
\par 听第 10 段材料，回答第 18 至 20 题，\par
\question{18. Why aren't more people willing to talk about work?}
\begin{choices}(3)
  \item Because they are fully occupied.
  \item Because they become lazier than before.
  \item Because they think working stress is too heavy.
\end{choices}
\question{19. What kind of job is not mentioned as an occupation?}
\begin{choices}(3)
  \item Engineer.
  \item Doctor.
  \item Cleaner.
\end{choices}
\question{20. What attitude should we take towards ordinary work?}
\begin{choices}(3)
  \item Indifferent.
  \item Respectful.
  \item Sympathetic.
\end{choices}
\parttitle{第二部分 阅读理解（共两节，满分 50 分）}
\sectitle{第一节（共15小题：每小题2.5分，满分37.5分）}
\par 阅读下列短文，从每题所给的 A、B、C 和 D 四个选项中，选出最佳答案，\par
\passagelabel{A}
\par The social purpose of the copyright system is to encourage creativity and the advancement of knowledge by giving those who make intellectual(智力的) contributions "unique right to their writings and discoveries" for limited times. Protection is provided for a much longer time for copyrights than for patents. The U.S. Copyright Law of 1976 provides a term or the lifetime of the author plus 50 years. It recognizes that the value of written works often extends over a longer time than that of technological advances in the public area.\par
\par It is essential to register for formal copyright protection in order to have some degree of protection against infringement(侵权) of one's works by others. The procedure for obtaining a copyright registration is quite simple—an application is submitted with a small fee and two copies of the work. No examination is made, and the copyright is immediately in force.\par
\par The Copyright Office will provide on request a variety of leaflets(小册子) describing the procedure and giving information about the classes of subject matter that can be registered.\par
\par (1) Books and newspapers (2) Lectures\par
\par (3) Dramatic compositions (4) Films\par
\par (5) Maps (6) Works of art\par
\par (7) Sound recordings (8) Photographs (9) Prints and labels (10) Musical compositions\par
\par Those are not subject to protection: (1) Ideas (2) Systems, methods and plans (3) Government publications (4) Laws and opinions of courts\par
\par You can contact the Copyright Office(www.copyrightfice.org) for more information.\par
\question{21. In which aspect do copyrights differ from patents?}
\begin{choices}(2)
  \item The application fee.
  \item The duration of protection.
  \item The intellectual value.
  \item The contribution to society
\end{choices}
\question{22. Which of the following should you provide for copyright registration?}
\begin{choices}(2)
  \item Copies of your work.
  \item Photos of yourself.
  \item Introductory leaflets.
  \item Examination certificates.
\end{choices}
\question{23. According to the text, the Copyright Office may reject \blanknum{}.}
\begin{choices}(2)
  \item a movie about a new method of time travel
  \item an original musical composition for children
  \item a book about one's opinions on American laws
  \item an official document issued by the local government
\end{choices}
\passagelabel{B}
\par The path to success isn't always clear-cut. Julia Green knows this well. Just as she graduated from college, a financial crisis swept the country, making it hard to get a job in journalism. Despite her growing worries, she didn't give up on her dream; instead, she tried a global volunteer program for her self-growth.\par
\par A year later, Green began taking random writing jobs to make ends meet, while dreaming of a day when she could write freely and be paid for it. Afterwards, she applied for the fully funded Master of Fine Arts program, and was accepted to study fiction writing at the University of Michigan. Though the program didn't end up opening many doors for her, she was paid to write and take classes.\par
\par Green has never shied away from challenging topics. Marriage, healthcare, employment, and more have featured noticeably in her writing, and she is especially impressive when it comes to music criticism. "I’ve intentionally tried to keep my lane(道路) very wide" she said. "It was what I had to do in the beginning. I wanted to become a well-rounded writer." She penned articles for a variety of websites. Slowly, she began to build a reputation for fearlessness.\par
\par "I've been writing since I was 22. This craft requires that you have a tough skin." said Green, when asked about her writing journey. "As a writer, you are not well served by caring what other people think about you. You have to be emotionally prepared to face rejection."\par
\par In the past decade, Green, lacking many examples to model after, has achieved remarkable success through great effort. She landed regular writing positions, had worked in respected journals, made the Forbes 30 Under 30 list, and wrote a world-class bestseller. By any young journalist's definition, Green had arrived.\par
\question{24. What do we know about Green from the first two paragraphs?}
\begin{choices}(2)
  \item A financial crisis ruined her dream.
  \item Further education made her optimistic.
  \item Volunteering brought her new jobs.
  \item Her passion for writing did not change.
\end{choices}
\question{25. Why did Green pursue a diverse path?}
\begin{choices}(2)
  \item To show off her creativity.
  \item To please critical readers.
  \item To develop her adaptability.
  \item To build her a reputation.
\end{choices}
\question{26. What matters most to her career according to Green?}
\begin{choices}(2)
  \item Handling criticism bravely.
  \item Maintaining youthful skin.
  \item Gaining public recognition.
  \item Writing about her journey.
\end{choices}
\question{27. Which words can best describe Green?}
\begin{choices}(2)
  \item Independent and considerate.
  \item Generous and creative.
  \item Hard-working and determined.
  \item Ambitious and modest.
\end{choices}
\passagelabel{C}
\par If someone asked whether you like the arts, you'd probably say you do at least in theory. According to a survey, more than two-thirds of U.S. adults say the arts "lift me up beyond everyday experiences". However, only 30 percent attended a concert of any type in 2017; 23 percent went to an art museum. Fewer than half actively created art of any kind.\par
\par The mismatch can boil down to the fact that we are weighed down by our day-to-day responsibilities, leaving our schedule packed. Maybe you like to play a little background music while you work or do the chores, but even before the pandemic, most of us rarely, if ever, saw a live performance, let alone visited a gallery or watched a play.\par
\par Too often, we let the dull reality of life get in the way of the arts. But this is a mistake. The arts are the opposite of an escape from reality; they might just be the most realistic glimpse we ever get into the nature and meaning of life. If you make time for consuming and producing the arts the same way you make time for work and exercise and family commitments, I assure you that you'll find your life getting fuller and happier. Think of a time when you heard a piece of music and wanted to cry. Or maybe your dizziness struck as you emerged from a narrow side street in an unfamiliar city. There you found yourself in a beautiful town square as if in a fantasy. They probably stimulated a sudden awakening, much like the shock from a lungful of pure oxygen after breathing in smoggy air.\par
\par If you are among the people who feel that the arts are pure pleasures to experience and participate in, you might see them as a luxury item, while a preferable attitude is to treat the arts less like a distracting pleasure, and more like exercise or sleep, a necessity. Then draw up a schedule of your art exposure journey, and gradually weave the arts into your everyday life.\par
\question{28. What prevents people from enjoying the arts according to the text?}
\begin{choices}(2)
  \item Lack of interest.
  \item Avoidance of duty.
  \item Shortage of time.
  \item Art space inaccessibility.
\end{choices}
\question{29. What can we learn about the role of the arts according to the text?}
\begin{choices}(2)
  \item They explore mysteries of life
  \item They spark emotional responses.
  \item They enhance physical well-being.
  \item They offer a getaway from daily life.
\end{choices}
\question{30. What will probably be talked about in the following paragraph?}
\begin{choices}(2)
  \item Barriers to combine the arts and life
  \item Benefits of engaging with the arts
  \item Forms of distracting pleasures.
  \item Ways to integrate the arts into daily routine.
\end{choices}
\question{31. Which is the most suitable title for the text?}
\begin{choices}(2)
  \item The Arts Journey: Refreshing Soul
  \item The Arts Pursuit: Transforming Dull into Full
  \item Embracing the Arts: From Luxury to Necessity
  \item Connecting with the Arts: From Reality to Fantasy
\end{choices}
\passagelabel{D}
\par Glasgow University researchers have created a laser(激光) camera using AI and quantum technologies to remotely detect a person's heartbeat, potentially changing health monitoring. Professor Faccio imagines various applications, such as setting up this technology in shopping malls or fixing laser heart monitors in homes. These monitors could be part of an extensive system tracking different health aspects, such as blood pressure and slight changes in movement, indicating early signs of conditions like Alzheimer's.\par
\par Monitoring a person's heartbeat from a distance holds significant value, as it can detect irregularities like abnormal heart rates, providing crucial warnings of potential heart issues. In contrast to the traditional use of stethoscopes(听诊器), this system overcomes limitations associated with proper stethoscope usage, such as the need for training and the potential reduction of heartbeat signals.\par
\par The system developed by Faccio and his team involves high-speed cameras which can record images at speeds of 2,000 frames per second. A laser beam is shone onto the skin of a person's throat and the reflections are used to measure exactly how much their skin is rising and falling as their main artery(动脉) expands and contracts as blood is forced through it. These changes involve movements of only a few billionths of a meter.\par
\par Such acuity is striking, though on its own the tracking of these tiny changes would not be enough to track a heartbeat. Other, much larger movements occur on a person's chest—from their breathing, for example—which would flood signals from their heartbeat.\par
\par "That is where AI comes in," Faccio said. "We use advanced computing systems to separate out everything except the vibrations(振动) caused by a person’s heartbeat. We know the frequency range of the human heartbeat, and AI focuses on that. The prime use of this technology, which we hope to have ready next year, will be to measure heartbeats easily and quickly outside hospitals. The benefits could be considerable."\par
\question{32. What can the laser camera be used for?}
\begin{choices}(2)
  \item Improving the usage of the stethoscopes.
  \item Measuring blood pressure remotely.
  \item Recording images of the human throat.
  \item Detecting a person's heartbeat from a distance.
\end{choices}
\question{33. What does paragraph 3 focus on as for the system?}
\begin{choices}(2)
  \item How it works.
  \item Where it is used.
  \item Whom it benefits.
  \item What it is composed of.
\end{choices}
\question{34. Which is closest in meaning to "acuity" in paragraph 4?}
\begin{choices}(2)
  \item Sensitivity.
  \item Complexity.
  \item Consistency.
  \item Intensity.
\end{choices}
\question{35. What role does AI play in the system?}
\begin{choices}(2)
  \item To measure skin movements.
  \item To recognize heartbeat vibrations.
  \item To record images at high speeds.
  \item To shine a laser beam onto the skin.
\end{choices}
\sectitle{第二节（共10小题；每小题1.5分，满分15分）}
\par 阅读下面短文，从短文后的选项中选出能填入空白处的最佳选项。选项中有两项为多余选项。\par
\par I am a child of traditional book reading. I used to go through one book continuously from beginning to end, taking notes or highlighting lines before going on to the next book. It was a slow but deeper way to absorb a book. \blanknum{36} My struggle with Dostoevsky's The Brothers Karamazov early on is a good example.\par
\par After college, my job involved reading a lot of background materials for projects. Many nights I stayed awake till past midnight leafing through pages. \blanknum{37} It helped me breeze through my reading requirements. However, for books I was genuinely passionate about, I switched to a slower pace, and still read one book at a time.\par
\par \blanknum{38} I no longer stick to one book, but jump between them, spending about half an hour on each. Currently, I'm reading six books in various stages, with more lined up to start soon. My speed reading skills help me scan quickly without losing comprehension. I can easily switch to slow reading mode whenever needed.\par
\par Recently, I learnt from a journal that my previous reading method is known as "Blocked Practice", which involves studying one topic thoroughly before moving on to another. \blanknum{39} However, it tends to restrict how you\par
\par apply knowledge in new and varied situations\par
\par On the other hand, the "mixing it up" approach, also referred to as "Interleaved Practice", could produce dramatic, long-lasting learning benefits. By switching between topics, readers can make connections between ideas and skills. \blanknum{40} Over time, readers could be more capable of solving problems critically. If this is true, then my evolution as a reader of books has been in the right direction all along.\par
\begin{choices}(1)
  \item It helps widen perspectives.
  \item This method has its limitations.
  \item That was when I took up speed reading.
  \item I'm not great at trying out different reading styles.
  \item That's why it took me months to finish heavy books.
  \item Nowadays I've acquired the habit of multiple book reading.
  \item This approach can be effective for mastering individual concepts.
\end{choices}
\parttitle{第三部分 语言知识运用（共两节，满分 30 分）}
\sectitle{第一节（共15小题；每小题1分，满分15分）}
\par 阅读下面短文，从短文后各题所给的 A、B、C、D 四个选项中，选出可以填入空白处的最佳选项。\par
\par My father's thrift(节俭) philosophy shaped my youth. Unlike peers receiving pocket money, I never got a penny not from cruelty, but to build \blanknum{41}. Each morning, he'd note, "Mrs. Wilson needs her garden weeded; Mr. Carter requires help \blanknum{42} bottles." I earned coins through \blanknum{43} work in the neighborhood: sweeping leaves, shoveling snow and planting trees.\par
\par Savings filled the tin under my bed, its weight showing my growing \blanknum{44}. Though he occasionally borrowed for emergencies(a flat tire or a medical bill), \blanknum{45} was exact and timely. When I lost two months' earnings in a(n) \blanknum{46} lemonade business, his only words were, "Live and learn". At twelve, I considered this cruel; at thirty, I \blanknum{47} its wisdom—true lessons grow through doing.\par
\par Modern thrift education \blanknum{48} saving, but ignores its key point: earning. Father believed "earned coins teach value; given coins generate greed" —a truth wealthy families overlook while \blanknum{49} thrift. By delivering papers at dawn and mowing lawns at dusk, I learned that true thrift is a triad(三位一体)—earn through \blanknum{50}, save with intention, spend with care. This trio, \blanknum{51} in classrooms, roots in blistered(起水泡的) hands and difficult situations.\par
\par \blanknum{52}, supermarkets throw outdated bread(still eatable) as shelters beg for donations. Witnessing global contrasts inspires my resolve: 1.3 billion tons of food go bad annually—enough to feed all the hungry—while 800 million starve. This phenomenon \blanknum{53} my Food Rescue project—redirecting bakery remains to food banks and training youth to \blanknum{54} deserted furniture for schools. Through such acts, thrift becomes social justice. As Father said, "Earning teaches respect for resources." In our wasteful era, thrift transforms from \blanknum{55} discipline to collective duty.\par
\cz{41}{self-control}{self-reliance}{self-improvement}{self-defense}
\cz{42}{breaking}{making}{delivering}{sorting}
\cz{43}{seasonal}{permanent}{dangerous}{voluntary}
\cz{44}{doubt}{fear}{pride}{respect}
\cz{45}{punishment}{repayment}{investment}{donation}
\cz{46}{failed}{expanding}{illegal}{profitable}
\cz{47}{ignore}{forget}{question}{value}
\cz{48}{gives up on}{focuses on}{benefits from}{moves beyond}
\cz{49}{opposing}{neglecting}{advocating}{realizing}
\cz{50}{labor}{creativity}{skills}{talents}
\cz{51}{understandable}{unteachable}{popular}{boring}
\cz{52}{Accidentally}{Occasionally}{Locally}{Temporarily}
\cz{53}{fuels}{delays}{replaces}{complicates}
\cz{54}{regain}{remake}{repair}{reuse}
\cz{55}{strict}{strong}{personal}{internal}
\sectitle{第二节（共10小题；每小题1.5分，满分15分）}
\par 阅读下面短文，在空白处填入 1 个适当的单词或括号内单词的正确形式.\par
\par The U. S. National Museum of Asian Art has returned two volumes of ancient Chinese silk manuscript(手稿) to China, which arrived in Beijing on Sunday \blanknum{} \blanknum{56}\blanknum{} nearly 80 years abroad.\par
\par This manuscript, whose value experts say is \blanknum{57} (compare) to that of the Dead Sea Scrolls in the West, is the oldest Chinese silk writing found so far and the only \blanknum{58} (know) one from the Warring States Period found in China. The name, Zidanku Silk Manuscript, came from the site \blanknum{59} it was stolen by grave robbers. It was then taken illegally to the U.S. in 1946.\par
\par China's ambassador to the U.S. said that the manuscript has over 900 characters \blanknum{60} is full of myths and numbers. He said \blanknum{61} return of the two volumes is a cultural milestone and a testament(证明) to China—U.S. cooperation in heritage protection and that the move \blanknum{62} (bring) new vitality to people-to-people exchanges and friendship. He also added that when the two countries work together on the basis of \blanknum{63} (equal) and respecting each other, they can achieve important results that help both sides and others.\par
\par An expert in international law mentioned that it is difficult to get back China's lost cultural items overseas because many of them \blanknum{64} (lose) over a long time for different reasons, and they are spread out. He said it is essential \blanknum{65} (make) a full list of these lost items and study where they came from and how they were moved.\par
\parttitle{第四部分 写作（共两节，满分 40 分）}
\sectitle{第一节（满分15分）}
\par 你校新来的外教 David 因故未能参加学校开学典礼，但在新老师介绍环节进行了微信视频通话。请你就此事给学校英文广播站投稿，内容包括：\par
\par (1) 有趣的视频通话过程;\par
\par (2) 你的期望和祝福，注意：\par
\par (1) 写作词数应为 80 左右;\par
\par (2) 可适当增加细节，以使行文连贯，\par
\par We Chat in English\par
\sectitle{第二节（满分25分）}
\par 阅读下面材料，根据其内容和所给段落开头语续写两段，使之构成一篇完整的短文。\par
\par In a house far from the city lived John and his mother. John's only hobby was wandering in the nearby picturesque forest. In its center was a shiny clear lake with schools of fish dancing in it. Lovely animals such as deer or rabbits were playing around and cute birds were singing in trees.\par
\par Fascinated by the beauty and liveliness of nature, he desired to become a great artist and present all the breath-taking scenes to the world.\par
\par John studied in the local primary school. Life was not easy for John and his mom after his father died. Knowing what hard work his mother did to make ends meet, John really didn't want his mother to worry. His mother knew the significance of education in life but she was against his dream of becoming an artist. "The profession of artist can't aid a person in meeting his daily expenses," she would say. But as each day passed, John became more passionate about his dream.\par
\par Due to the living circumstances, even buying some new colors would be a luxury. John only used the worn drawing file and limited colors for his daily practice. Though lacking fine painting materials, he never gave up drawing. He spent a significant amount of time staying in the forest, secretly perfecting his craft. With continuous hard work and a distinct gift for art, he drew beautifully.\par
\par One day John read about a competition being held in the town, whose deadline was only two days away. He got the perfect opportunity to fulfill his desire but he was in a dilemma(两难境地) as for the competition he needed special paints and some brushes. He didn't have the courage to ask his mother for money as he knew her prejudice against art.\par
\par That evening his mom returned from work and handed him a few paper notes, saying, "This is for your lunch at school next month." The money was enough for his desired materials! John's face lit up. 注意:\par
\par (1) 续写词数应为 150 左右;\par
\par (2) 请按如下格式在答题卡的相应位置作答。\par
\writingpara{Paragraph 1: Early the next morning, John secretly left home with the money. \blanknum{}}
\writinglines{8}
\writingpara{Paragraph 2: A week later; a phone call came to John's home.}
\writinglines{8}
\end{document}
